\section{多项式与生成函数}

\subsection{多项式卷积 (Polynomial Convolution)}

\subsubsection{快速傅里叶变换 (FFT)}

\textbf{用途：}
在 $O(n\log n)$ 时间内计算两个多项式的乘积。给定两个多项式 $A(x)=\sum_{i=0}^{n-1}a_ix^i$ 和 $B(x)=\sum_{i=0}^{m-1}b_ix^i$，FFT 可以高效求出它们的卷积 $C(x)=A(x)\cdot B(x)$ 的系数。

\textbf{核心思想：}
\begin{itemize}
    \item \textbf{点值表示：} 一个 $n$ 次多项式由它在 $n$ 个不同点处的取值唯一确定。
    \item \textbf{单位根：} 选择 $\omega_n^k = e^{2\pi i k/n}$ 作为求值点，利用其周期性和对称性进行分治。
    \item \textbf{蝴蝶变换：} 在迭代实现中，将下标按二进制位逆序重排，使得可以直接从底层向上合并，避免递归。
    \item \textbf{逆变换：} 只需对点值再做一次 FFT，并将结果除以 $n$（或者将单位根替换为共轭后对结果取共轭并除以 $n$）。
\end{itemize}

\textbf{算法流程（迭代版）：}
\begin{enumerate}
    \item 将多项式 $A,B$ 补全至长度 $N=2^k$，且 $N$ 不小于乘积的次数。
    \item 进行蝴蝶重排，将序列按二进制逆序重新排列。
    \item 枚举合并步长 $\text{len}=2,4,\dots,N$，对应单位根 $\omega_{\text{len}}=e^{\pm 2\pi i/\text{len}}$，合并相邻的两段。
    \item 完成正向 FFT 后得到点值，对应点值相乘，再执行一次逆 FFT 得到卷积结果（注意除以 $N$ 或缩放）。
\end{enumerate}

\textbf{时间复杂度：} 蝴蝶重排 $O(N)$，逐层合并共 $\log N$ 层，每层 $O(N)$，总复杂度 $O(N\log N)$。空间复杂度 $O(N)$。

\lstinputlisting[language=C++]{codes/08-polynomial-01.cpp}

\textbf{说明：}
\begin{itemize}
    \item 代码中使用了 \texttt{std::complex<double>}，实际应用中若需避免浮点误差，可改用整数意义下的 NTT（数论变换）。
    \item 最终系数存储在 \texttt{double} 向量中，对于整数多项式可进行四舍五入。
\end{itemize}

\subsubsection{快速数论变换 (NTT)}

\textbf{用途：}
在 $O(n\log n)$ 时间内计算两个\textbf{整数系数}多项式在模意义下的乘积，避免浮点误差。
给定多项式 $A(x),B(x)$ 和一个 NTT-friendly 的质数模数 $P$（如 $998244353$，其原根为 $3$），NTT 能精确计算 $C(x)=A(x)\cdot B(x) \bmod P$。

\textbf{核心思想：}
\begin{itemize}
    \item 类比 FFT 中的复单位根 $\omega_n$，在模 $P$ 下使用\textbf{原根} $g$ 构造\textbf{主单位根} $\omega_n = g^{(P-1)/n}$，满足完全相同的代数性质。
    \item 需要 $n \mid (P-1)$，即 $P$ 形如 $k\cdot 2^m+1$，保证变换长度 $n$ 为 $2$ 的幂次。常用 $P=998244353=119\times 2^{23}+1$，原根 $g=3$。
    \item 蝴蝶变换与 FFT 一致，逆变换时使用单位根的逆元并在最后乘 $n^{-1}\bmod P$。
\end{itemize}

\textbf{原根与单位根的对应关系：}

NTT 能够在模 $P$ 下完美复现 FFT 的所有性质，根本原因在于\textbf{原根} $g$ 的阶为 $P-1$，令 $\omega_n = g^{(P-1)/n}$，则 $\omega_n$ 恰为模 $P$ 下的 $n$ 次\textbf{主单位根}，其代数性质与复数域中的 $\omega_n = e^{2\pi i/n}$ 一一对应。

\begin{tabular}{|c|c|c|}
\hline
\textbf{性质} & \textbf{复数域 (FFT)} & \textbf{模域 (NTT)} \\ \hline
主单位根定义 & $\omega_n = e^{2\pi i/n}$ & $\omega_n = g^{(P-1)/n} \pmod{P}$ \\
\hline
$\omega_n^n = 1$ & $e^{2\pi i}=1$ & $\left(g^{(P-1)/n}\right)^n = g^{P-1} = 1$ \\
   & 旋转一周回到原处 & 费马小定理 $g^{P-1}\equiv 1$ \\
\hline
$\omega_n^{n/2} = -1$ & $e^{\pi i} = -1$ & $g^{(P-1)/2} \equiv -1$ \\
   & 半周翻转 & 二次剩余意义下的 $-1$ \\
\hline
周期性 & $\omega_n^k = \omega_n^{(k \bmod n)}$ & $\omega_n^k = \omega_n^{(k \bmod n)}$ \\
   & 以 $n$ 为周期的旋转 & 指数模 $n$ 周期 \\
\hline
折半性 & $\omega_{2n}^{2k} = \omega_n^k$ & $\omega_{2n}^{2k} = g^{\frac{P-1}{2n}\cdot 2k} = g^{\frac{P-1}{n}\cdot k} = \omega_n^k$ \\
   & 加倍长度后偶次项不变 & 指数约分后等价 \\
\hline
\end{tabular}

\vspace{0.3cm}

\textbf{算法流程（迭代版）：}
\begin{enumerate}
    \item 将多项式补全至长度 $N=2^k$，满足 $N$ 不小于乘积的次数，且 $N \mid (P-1)$。
    \item 蝴蝶重排（二进制逆序）。
    \item 枚举步长，计算当前层的单位根 $w_{\text{len}}=g^{(P-1)/\text{len}}$，合并相邻块。
    \item 正变换得到点值，逐点相乘，逆变换后各乘 $N^{-1}\bmod P$ 得到最终系数。
\end{enumerate}

\textbf{时间复杂度：} 与 FFT 相同，为 $O(N\log N)$，全部使用整数运算，没有浮点误差。

\lstinputlisting[language=C++]{codes/08-polynomial-02.cpp}

\textbf{说明：}
\begin{itemize}
    \item 代码固定模数 $998244353$ 和原根 $3$，这是算法竞赛中最常用的 NTT 模数。
    \item 若需要其他模数，需保证 $P-1$ 含有足够大的 $2$ 的幂次，并找到对应的原根。
    \item 逆变换时单位根取逆，最后乘 $n^{-1}$ 即可得到卷积系数。
\end{itemize}

\subsubsection{无位逆序 NTT (Bit-reversal-free NTT)}

\textbf{用途：}
与上一节的标准 NTT 相同，在 $O(N\log N)$ 时间内计算模意义下的多项式卷积，但整个流程\textbf{不执行任何位逆序置换 (bit-reversal permutation)}：省去 $O(N)$ 的跨步交换、\code{rev} 数组与缓存不友好的随机访存，代码更短，实际常数更优。

\textbf{核心思想：}
\begin{itemize}
    \item 蝴蝶操作有两种方向，\textbf{均无需预先重排}。\textbf{按频率抽取 (DIF, Decimation-In-Frequency)}：输入自然序，输出自动变为位逆序；\textbf{按时间抽取 (DIT, Decimation-In-Time)}：输入位逆序，输出自动恢复自然序。
    \item DIF 与 DIT 互为转置，蝴蝶形状不同：DIF \textbf{先加减、后乘旋转因子}，得到 $(u+v,\ (u-v)\,\omega^{k})$，层长从大到小枚举；DIT \textbf{先乘旋转因子、后加减}，得到 $(u+v\omega^{k},\ u-v\omega^{k})$，层长从小到大枚举。
    \item 卷积的逐点相乘与下标顺序无关：记 $R$ 为位逆序置换矩阵，$F$ 为 DFT 矩阵，则 DIF 实现的是 $R\cdot F$；而置换是逐下标作用，$(Rx)\odot(Ry)=R(x\odot y)$，位逆序的点值逐点相乘后\textbf{仍是位逆序}。
    \item 逆 DIT 恰好实现 $F^{-1}\cdot R$（输入位逆序、输出自然序）。于是「DIF 正变换 $\to$ 逐点相乘 $\to$ DIT 逆变换」完美衔接：\[F^{-1}R\big((RFA)\odot(RFB)\big)=F^{-1}\big((FA)\odot(FB)\big)=A\cdot B,\] 位逆序状态只存在于中间，任何显式置换都不需要。
\end{itemize}

\textbf{算法流程：}
\begin{enumerate}
    \item 将 $A,B$ 补零至长度 $N=2^k$（$N\ge n+m-1$，且 $N\mid P-1$）。
    \item 对 $A,B$ 分别执行 DIF：枚举块长 $\text{len}=N,N/2,\dots,2$，层单位根 $\omega_{\text{len}}=g^{(P-1)/\text{len}}$，块内前半放 $u+v$，后半放 $(u-v)\,\omega_{\text{len}}^{k}$，得到位逆序的点值。
    \item 逐点相乘 $C_i=A_iB_i\bmod P$（两序列下标顺序一致，直接对应相乘即可）。
    \item 对 $C$ 执行逆 DIT：枚举块长 $\text{len}=2,4,\dots,N$，旋转因子取逆元 $\omega_{\text{len}}^{-k}$，块内放 $u\pm v\,\omega_{\text{len}}^{-k}$，最后整体乘 $N^{-1}\bmod P$，即得自然序的卷积系数。
\end{enumerate}

\textbf{时间复杂度：} 仍为 $O(N\log N)$；相比标准实现省去位逆序置换的 $O(N)$ 次跨步交换与 \code{rev} 数组开销，实际运行常数更小。

\lstinputlisting[language=C++]{codes/08-polynomial-03.cpp}

\textbf{说明：}
\begin{itemize}
    \item 正逆变换必须「DIF 配 DIT」成对使用：两次都用 DIF 或两次都用 DIT 都得不到正确的卷积。
    \item DIF 的输出与逆 DIT 的输入均为位逆序；若后续算法需要\textbf{自然序点值}（如多项式多点求值、跨操作保留点值复用），仍需显式做一次位逆序置换。
    \item 两者唯一的实现区别是：旋转因子乘在减法\textbf{之后}（DIF）或\textbf{之前}（DIT），以及层长枚举方向相反（$N\to2$ 对 $2\to N$）。
    \item 适用前提是变换全程「正逆成对」出现：多项式乘法及基于卷积的模板（求逆 / 开方 / \code{ln} / \code{exp}）中的每次正变换都有对应逆变换且中间无需自然序点值时，可用本模板替换之前的 \code{ntt} 实现。
\end{itemize}

\subsubsection{卷积与相关 (Convolution and Correlation)}

\textbf{用途：}
FFT/NTT 不仅能计算多项式乘法，也能直接套用于\textbf{序列}的卷积 (Convolution) 与相关 (Correlation) 运算，在字符串匹配、模式匹配、计数问题等题目中经常出现。本节统一以模 $998244353$ 的 NTT 模板给出普通卷积、循环卷积、普通互（自）相关、循环互（自）相关四种运算的定义、实现与注意事项。

\begin{itemize}
    \item \textbf{普通卷积 (Linear Convolution)：} $c_k=\sum_{i+j=k}a_ib_j$（越界下标视为 $0$），等价于多项式乘法。实现时必须将 $a,b$ 补零至长度 $N\ge n+m-1$（$N$ 取最小 $2$ 的幂）再变换；否则长度不足时，超出部分会\textbf{回绕 (wrap-around)} 混叠到结果头部，得到的是循环卷积而非普通卷积。
    \item \textbf{循环卷积 (Circular Convolution)：} $c_k=\sum_{i+j\equiv k\pmod N}a_ib_j$，即长度为 $N$ 的周期卷积。当 $N\ge n+m-1$ 时与普通卷积完全一致。实现方式二选一：其一，$N$ 为 $2$ 的幂且 $N\mid(P-1)$ 时，直接做长度 $N$ 的 NTT 后点值相乘、逆变换即得；其二，任意 $N$ 时，先算普通卷积，再把结果按 $k\bmod N$ 折叠回绕相加。
    \item \textbf{普通互（自）相关 (Linear Cross-Correlation)：} $r_k=\sum_i a_ib_{i+k}$，$k\in[-(m-1),n-1]$，衡量两序列的相似程度。核心恒等式 $r_k=(a\ast\tilde b)[m-1-k]$，其中 $\tilde b_j=b_{m-1-j}$ 为翻转后的 $b$，即\textbf{翻转 $b$ 后做普通卷积、再镜像回原顺序}。取 $b=a$ 即自相关：$r_{-k}=r_k$（偶对称），$r_0=\sum_i a_i^2$ 为序列\textbf{能量}。
    \item \textbf{循环互（自）相关 (Circular Cross-Correlation)：} $r_k=\sum_i a_ib_{(i+k)\bmod N}$，$k\in[0,N-1]$。实现：翻转 $b$ 后做长度为 $N$ 的循环卷积得 $c$，再由 $r_k=c_{(m-1-k)\bmod N}$ 镜像。取 $b=a$ 得循环自相关，满足 $r_k=r_{(N-k)\bmod N}$，可用于检测序列周期性。
\end{itemize}

\textbf{注意事项：}
\begin{itemize}
    \item \textbf{补零长度：} 普通卷积要求 $N\ge n+m-1$，否则回绕混叠；$N$ 取不小于 $n+m-1$ 的最小 $2$ 的幂。
    \item \textbf{NTT 长度限制：} 模 $998244353$ 要求 $N\mid P-1$，即 $N\le 2^{23}$；数据规模更大时需多模 NTT + CRT 或改用 FFT。
    \item \textbf{精度与模数：} FFT 为浮点运算，整数结果需四舍五入取整；NTT 全程整数无误差，但结果为 $\bmod P$ 意义下，真实值可能超过 $P$ 时需多模 NTT。
    \item \textbf{下标偏移：} 相关结果的下标与卷积不同：\code{cross\_correlation} 返回 $res[k+(m-1)]=r_k$，零移位对应 $res[m-1]$，取答案时勿混淆。
    \item \textbf{共轭与翻转：} 频域相关定理 $R=\overline{\operatorname{DFT}(A)}\cdot\operatorname{DFT}(B)$ 仅适用于复数 FFT；NTT 模域没有共轭运算，一律用「翻转 $b$ 后卷积」实现，两者在代数上等价。
\end{itemize}

\lstinputlisting[language=C++]{codes/08-polynomial-04.cpp}

\subsection{生成函数 (Generating Functions)}

\subsubsection{普通生成函数与指数生成函数 (OGF and EGF)}

\textbf{用途：}
生成函数 (母函数) 将无穷数列 $\{a_n\}_{n\ge 0}$ 编码为形式幂级数，把计数与递推问题转化为幂级数之间的代数运算，再配合本节后面的多项式模板（求逆 / 开方 / \code{ln} / \code{exp}）在 $O(n\log n)$ 时间内求出数列的通项或前 $n$ 项。

\textbf{核心概念：}
\begin{itemize}
    \item \textbf{普通生成函数 (OGF, Ordinary Generating Function)：} $F(x)=\sum_{n\ge 0}a_nx^n$。两个 OGF 相乘对应序列的\textbf{卷积}：$c_n=\sum_{i+j=n}a_ib_j$，组合意义为「无序拼接」——把问题拆成两个独立子问题再合并方案数。
    \item \textbf{指数生成函数 (EGF, Exponential Generating Function)：} $\hat F(x)=\sum_{n\ge 0}a_n\dfrac{x^n}{n!}$。两个 EGF 相乘时 $c_n=\sum_{i}\binom{n}{i}a_ib_{n-i}$，组合意义为「标号拆分」——自动携带把 $n$ 个标号元素分配给两部分的方案数，适合排列、集合划分等\textbf{带标号}计数。
    \item \textbf{形式幂级数：} 生成函数只关心系数、不考虑收敛性，加法、乘法、求导、复合等规则与多项式完全一致。
    \item \textbf{常用运算对应：} 乘 $x$ 对应下标整体 $+1$；求导 $F'$ 对应取出 $a_{n+1}$；$\dfrac{F(x)}{1-x}$ 对应前缀和；$\bmod\ x^n$ 对应「只保留前 $n$ 项」（多项式模板的语义）。
    \item \textbf{建模套路：} 对递推式两端建立生成函数，转化为代数方程求解。例如 Catalan 数 $C_{n+1}=\sum_{i}C_iC_{n-i}$ 解得 $C(x)=\dfrac{1-\sqrt{1-4x}}{2x}$，直接用多项式开方 + 求逆即可批量求值。
\end{itemize}

\textbf{常见数列的生成函数：}

\begin{center}
\begin{tabularx}{\textwidth}{|c|c|X|}
\hline
\textbf{数列 $a_n$} & \textbf{生成函数} & \textbf{说明} \\ \hline
$1,1,1,\dots$ & $\dfrac{1}{1-x}$ & 基础几何级数 (OGF) \\ \hline
$1,q,q^2,\dots$ & $\dfrac{1}{1-qx}$ & 等比数列 (OGF) \\ \hline
$\binom{m}{n}$（固定 $m$） & $(1+x)^m$ & 二项式定理 (OGF) \\ \hline
$\binom{n+m-1}{n}$ & $\dfrac{1}{(1-x)^m}$ & $n$ 次可重选取 (OGF) \\ \hline
$0,1,1,2,3,\dots$ (Fibonacci) & $\dfrac{x}{1-x-x^2}$ & $F_n=F_{n-1}+F_{n-2}$ (OGF) \\ \hline
$\dfrac{1}{n+1}\binom{2n}{n}$ (Catalan) & $\dfrac{1-\sqrt{1-4x}}{2x}$ & $C_{n+1}=\sum_{i=0}^{n}C_iC_{n-i}$ (OGF) \\ \hline
$n!$ & $\dfrac{1}{1-x}$ & 排列数 (EGF) \\ \hline
$S(n,k)$（固定 $k$） & $\dfrac{(e^x-1)^k}{k!}$ & 第二类斯特林数：$n$ 个标号元素分为 $k$ 个非空无序集合 (EGF) \\ \hline
贝尔数 $B_n$ & $e^{e^x-1}$ & 集合划分总数 $\sum_k S(n,k)$ (EGF) \\ \hline
错位排列 $D_n$ & $\dfrac{e^{-x}}{1-x}$ & $D_n=n!\sum_{i=0}^{n}\dfrac{(-1)^i}{i!}$ (EGF) \\ \hline
\end{tabularx}
\end{center}

\textbf{说明：}
\begin{itemize}
    \item OGF 适合无标号 / 序列型组合，EGF 适合标号组合；同一数列的 OGF 与 EGF 一般不同，查表时注意区分。
    \item 表中 $e^x=\sum_{n\ge 0}\dfrac{x^n}{n!}$ 可视为「单个标号元素集合」的 EGF，是构造斯特林数、贝尔数等 EGF 的基本积木。
\end{itemize}

\subsubsection{多项式求逆 (Polynomial Inverse)}

\textbf{用途：}
给定 $n$ 项多项式 $A(x)$（$a_0\not\equiv 0$），求 $B(x)$ 的前 $n$ 项使得
\[ A(x)\,B(x) \equiv 1 \pmod{x^n}. \]
多项式求逆是开方、\code{ln}/\code{exp}、除法取模的\textbf{公共底层依赖}，也是生成函数中「$1/F(x)$」这一除法运算的直接实现。

\textbf{核心思想：}
\begin{itemize}
    \item \textbf{Newton 迭代（倍增）：} 已知 $B_t \equiv A^{-1} \pmod{x^t}$，则
    \[ B_{2t} \;=\; B_t\,\bigl(2 - A\,B_t\bigr) \pmod{x^{2t}}, \]
    误差以平方速度缩减，精度每轮翻倍。
    \item \textbf{初值：} $B_1 = a_0^{-1} \bmod p$，单点逆元用费马小定理求得。
    \item \textbf{NTT 长度：} 每轮变换长度取 $2k$，保证 $A B_t^2$ 的次数界不超过 $2k-3$，不会发生循环卷积回绕。
\end{itemize}

\textbf{时间复杂度：} 递推 $T(n)=T(n/2)+O(n\log n)=O(n\log n)$。空间复杂度 $O(n)$。

\lstinputlisting[language=C++]{codes/08-polynomial-05.cpp}

\textbf{说明：}
\begin{itemize}
    \item 代码自包含（内嵌与 08-polynomial-02 相同的 NTT），可整文件复制到比赛工程中使用。
    \item 若 $a_0 \equiv 0$ 则逆不存在，题面一般保证 $a_0 \ne 0$。
    \item 生成函数建模中「$\bmod\ x^n$」正对应「只看前 $n$ 项」，多项式求逆即该语义下的 $1/A$。
\end{itemize}

\subsubsection{多项式开方 (Polynomial Square Root)}

\textbf{用途：}
给定 $n$ 项多项式 $A(x)$，求 $B(x)$ 的前 $n$ 项使 $B(x)^2 \equiv A(x) \pmod{x^n}$。生成函数中大量形如 $F=xF^2+1$（Catalan）、$F^2=F+x$ 的代数方程都归结为多项式开方。

\textbf{核心思想：}
\begin{itemize}
    \item \textbf{Newton 迭代：} 已知 $B_t^2 \equiv A \pmod{x^t}$，则
    \[ B_{2t} \;=\; \frac{1}{2}\left(B_t + A\,B_t^{-1}\right) \pmod{x^{2t}}, \]
    其中 $B_t^{-1}$ 用多项式求逆（上一小节）计算。
    \item \textbf{常数项：} 要求 $a_0$ 在模意义下可开方；本模板默认 $a_0=1$（竞赛题最常见），一般情形需先用 Cipolla 算法（二次剩余）求出 $\sqrt{a_0} \bmod p$ 作为初值。
\end{itemize}

\textbf{时间复杂度：} 每轮一次多项式求逆 + 一次乘法，$T(n)=T(n/2)+O(n\log n)=O(n\log n)$。

\lstinputlisting[language=C++]{codes/08-polynomial-06.cpp}

\subsubsection{多项式对数与指数 (Polynomial ln / exp)}

\textbf{用途：}
$\ln A$ 与 $\exp A=\sum_{k\ge 0}\dfrac{A^k}{k!}$ 是 EGF 运算的核心：多项式快速幂 $A^{p}=\exp(p\ln A)$（$p$ 可以是任意大的整数甚至有理数）、EGF 的组合构造（集合、序列构造）都依赖它们。

\textbf{核心思想：}
\begin{itemize}
    \item \textbf{对数（要求 $a_0=1$）：} 对 $\ln A$ 求导消去常数项后积分还原：
    \[ \ln A \;=\; \int \frac{A'}{A}\,\mathrm{d}x, \]
    即「求导 $\to$ 乘 $A^{-1} \to$ 积分」，积分时逐项 $b_i = \dfrac{c_{i-1}}{i}$。
    \item \textbf{指数（要求 $a_0=0$）：} 对 $G(E)=\ln E - A$ 做 Newton 迭代：
    \[ E_{2t} \;=\; E_t\,\bigl(1 + A - \ln E_t\bigr) \pmod{x^{2t}}. \]
    \item \textbf{多项式快速幂：} $A^p=e^{p\ln A}$，要求 $a_0=1$；若 $a_0\ne 1$，先提出 $x$ 的幂次与常数项单独处理。
\end{itemize}

\textbf{时间复杂度：} $\ln$ 为一次求逆 + 一次乘法 $O(n\log n)$；\code{exp} 每轮一次 $\ln$ + 一次乘法，同样 $O(n\log n)$。

\lstinputlisting[language=C++]{codes/08-polynomial-07.cpp}

\textbf{说明：}
\begin{itemize}
    \item $\ln$ 的常数项：$a_0\ne 1$ 时先整体除以 $a_0$ 再计算；$\ln a_0$ 在模意义下无初等求法，通常题目不涉及。
    \item \code{exp} 要求 $a_0=0$，这是 $e^{a_0}$ 在模意义下有定义的自然要求。
\end{itemize}

\subsubsection{多项式除法与取模 (Polynomial Division and Modulo)}

\textbf{用途：}
给定 $n$ 项多项式 $A$、$m$ 项多项式 $B$（$b_{m-1}\ne 0$），求唯一的 $Q,R$ 使
\[ A \;=\; Q\,B + R, \qquad \deg R < m-1. \]
典型应用：\textbf{常系数齐次线性递推}中把 $x^n$ 对特征多项式取模（Kitamasa / 多项式快速幂优化递推），以及多项式 GCD。

\textbf{核心思想：}
\begin{itemize}
    \item \textbf{系数反转：} 定义 $A^{R}(x)=x^{n-1}A(1/x)$（系数倒序排列）。由 $A=QB+R$ 且 $\deg R < m-1$ 可得
    \[ A^{R} \;\equiv\; Q^{R}\,B^{R} \pmod{x^{\,n-m+1}}, \]
    余式 $R$ 的次数过低，在 $\bmod\ x^{n-m+1}$ 意义下被完全消去。
    \item \textbf{求商：} $Q^{R} \equiv A^{R}\,(B^{R})^{-1} \pmod{x^{n-m+1}}$，用多项式求逆后相乘，再反转回 $Q$。
    \item \textbf{求余：} $R = A - Q B$，一次乘法后逐位相减。
\end{itemize}

\textbf{时间复杂度：} 两次多项式乘法 + 一次多项式求逆，$O(n\log n)$。

\lstinputlisting[language=C++]{codes/08-polynomial-08.cpp}
